% !TeX root = ../main.tex
% !TeX program = xelatex
% !TeX encoding = UTF-8
\section{结论}

本文围绕资源分配、设备适配、建筑动态与区域容量，建立源—网—荷—储统一模型及相互衔接的多目标、连续运行和物理反馈模块。共享资源预算约束直接利用与热泵提温的共同取热，精确RC状态与独立储热分别描述两类热量时移，共享电力上限把逐楼运行连接至区域可行性。

各机制实验验证了四项优势：固定实际日集合下，六栋外部计量建筑的归一分布误差中位数由12.3150\%降至4.4775\%；厂家Q/P有效域支持的三城六组同资产调度降低条件电量费1.1013\%--1.5970\%；纯热容二维对照量化了150.52元的建筑柔性—区域容量费用交互；有限网络反馈将最大压降梯度由1081.42降至237.96~Pa/m，并使最低远端供水达到57.45~$^\circ$C。

三城中点碳帽保留80.92\%--89.44\%的降费空间，条件投资给出随采购和维护预算变化的方案选择阈值。模型不仅给出候选方案，还能解释资源如何使用、调度为何改变及收益如何形成。模块化配置、伪代码对应的实现和独立物理账本共同支持不同城市背景下的复用与核验。
